\documentclass[12pt, titlepage]{article}

\usepackage{booktabs}
\usepackage{tabularx}
\usepackage{hyperref}
\hypersetup{
    colorlinks,
    citecolor=blue,
    filecolor=black,
    linkcolor=red,
    urlcolor=blue
}
\usepackage[round]{natbib}

\input{../Comments.text}
\input{../Common.text}

\begin{document}

\title{System Verification and Validation Plan for \projname{}}
\author{\authname}
\date{\today}
	
\maketitle

\pagenumbering{roman}

\section*{Revision History}

\begin{tabularx}{\textwidth}{p{3cm}p{2cm}X}
\toprule {\bf Date} & {\bf Version} & {\bf Notes}\\
\midrule
Date 1 & 1.0 & Notes\\
Date 2 & 1.1 & Notes\\
\bottomrule
\end{tabularx}

~\\
\wss{The intention of the VnV plan is to increase confidence in the software.
However, this does not mean listing every verification and validation technique
that has ever been devised.  The VnV plan should also be a \textbf{feasible}
plan. Execution of the plan should be possible with the time and team available.
If the full plan cannot be completed during the time available, it can either be
modified to ``fake it'', or a better solution is to add a section describing
what work has been completed and what work is still planned for the future.}

\wss{The VnV plan is typically started after the requirements stage, but before
the design stage.  This means that the sections related to unit testing cannot
initially be completed.  The sections will be filled in after the design stage
is complete.  the final version of the VnV plan should have all sections filled
in.}

\newpage

\tableofcontents

\listoftables
\wss{Remove this section if it isn't needed}

\listoffigures
\wss{Remove this section if it isn't needed}

\newpage

\section{Symbols, Abbreviations, and Acronyms}

\renewcommand{\arraystretch}{1.2}
\begin{tabular}{l l} 
  \toprule		
  \textbf{symbol} & \textbf{description}\\
  \midrule 
  T & Test\\
  \bottomrule
\end{tabular}\\

\wss{symbols, abbreviations, or acronyms --- you can simply reference the SRS
  \citep{SRS} tables, if appropriate}

\wss{Remove this section if it isn't needed}

\newpage

\pagenumbering{arabic}

This document ... \wss{provide an introductory blurb and roadmap of the
  Verification and Validation plan}

\section{General Information}

\subsection{Summary}

\wss{Say what software is being tested.  Give its name and a brief overview of
  its general functions.}

\subsection{Objectives}

\wss{State what is intended to be accomplished.  The objective will be around
  the qualities that are most important for your project.  You might have
  something like: ``build confidence in the software correctness,''
  ``demonstrate adequate usability.'' etc.  You won't list all of the qualities,
  just those that are most important.}

\wss{You should also list the objectives that are out of scope.  You don't have 
the resources to do everything, so what will you be leaving out.  For instance, 
if you are not going to verify the quality of usability, state this.  It is also 
worthwhile to justify why the objectives are left out.}

\wss{The objectives are important because they highlight that you are aware of 
limitations in your resources for verification and validation.  You can't do everything, 
so what are you going to prioritize?  As an example, if your system depends on an 
external library, you can explicitly state that you will assume that external library 
has already been verified by its implementation team.}

\subsection{Custom Documents}

\wss{Summarize the CDs (Custom Documents) that were tackled by this project.  CDs 
can include usability testing, code walkthroughs, user documentation, formal
proof, GenderMag personas, Design Thinking, etc.  CDs should have already
been approved by the course instructor as included in your problem statement.
You can use a pull request to update your CDs (in TeamComposition.csv or
Repos.csv) if your plan changes as a result of the VnV planning exercise.}

\subsection{Relevant Documentation}

\wss{Reference relevant documentation.  This will definitely include your SRS
  and your other project documents (design documents, like MG, MIS, etc).  You
  can include these even before they are written, since by the time the project
  is done, they will be written.  You can create BibTeX entries for your
  documents and within those entries include a hyperlink to the documents.}

\citet{SRS}

\wss{Don't just list the other documents.  You should explain why they are relevant and 
how they relate to your VnV efforts.}

\section{Plan}

\wss{Introduce this section.  You can provide a roadmap of the sections to
  come.}

\subsection{Verification and Validation Team}

\wss{Your teammates.  Maybe your supervisor.
  You should do more than list names.  You should say what each person's role is
  for the project's verification.  A table is a good way to summarize this information.}

\subsection{SRS Verification}

\wss{List any approaches you intend to use for SRS verification.  This may
  include ad hoc feedback from reviewers, like your classmates (like your
  primary reviewer), or you may plan for something more rigorous/systematic.}

\wss{If you have a supervisor for the project, you shouldn't just say they will
read over the SRS.  You should explain your structured approach to the review.
Will you have a meeting?  What will you present?  What questions will you ask?
Will you give them instructions for a task-based inspection?  Will you use your
issue tracker?}

\wss{Maybe create an SRS checklist?}

\subsection{Design Verification}

\wss{Plans for design verification}

\wss{The review will include reviews by your classmates}

\wss{Create a checklists?}

\subsection{Verification and Validation Plan Verification}

\wss{The verification and validation plan is an artifact that should also be
verified.  Techniques for this include review and mutation testing.}

\wss{The review will include reviews by your classmates}

\wss{Create a checklists?}

\subsection{Implementation Verification}

\wss{You should at least point to the tests listed in this document and the unit
  testing plan.}

\wss{In this section you would also give any details of any plans for static
  verification of the implementation.  Potential techniques include code
  walkthroughs, code inspection, static analyzers, etc.}

\wss{The final class presentation in CAS 741 could be used as a code
walkthrough.  There is also a possibility of using the final presentation (in
CAS741) for a partial usability survey.}

\subsection{Automated Testing and Verification Tools}

\wss{What tools are you using for automated testing.  Likely a unit testing
  framework and maybe a profiling tool, like ValGrind.  Other possible tools
  include a static analyzer, make, continuous integration tools, test coverage
  tools, etc.  Explain your plans for summarizing code coverage metrics.
  Linters are another important class of tools.  For the programming language
  you select, you should look at the available linters.  There may also be tools
  that verify that coding standards have been respected, like flake9 for
  Python.}

\wss{If you have already done this in the development plan, you can point to
that document.}

\wss{The details of this section will likely evolve as you get closer to the
  implementation.}

\subsection{Software Validation}

\wss{If there is any external data that can be used for validation, you should
  point to it here.  If there are no plans for validation, you should state that
  here.}

\wss{You might want to use review sessions with the stakeholder to check that
the requirements document captures the right requirements.  Maybe task based
inspection?}

\wss{For those capstone teams with an external supervisor, the Rev 0 demo should 
be used as an opportunity to validate the requirements.  You should plan on 
demonstrating your project to your supervisor shortly after the scheduled Rev 0 demo.  
The feedback from your supervisor will be very useful for improving your project.}

\wss{For teams without an external supervisor, user testing can serve the same purpose 
as a Rev 0 demo for the supervisor.}

\wss{This section might reference back to the SRS verification section.}

\section{System Tests}

\wss{There should be text between all headings, even if it is just a roadmap of
the contents of the subsections.}

\subsection{Tests for Functional Requirements}

\wss{Subsets of the tests may be in related, so this section is divided into
  different areas.  If there are no identifiable subsets for the tests, this
  level of document structure can be removed.}

\wss{Include a blurb here to explain why the subsections below
  cover the requirements.  References to the SRS would be good here.}
The tests below verify the functional requirements defined in Section 6 in the SRS. Tests are group
by major functionalitys including project proposals, team formation, project applications, approvals,
communication, and course administration. Each test evaluates whether the system produces the expected
results under the given conditions

\subsubsection{Project Proposal Submission and Validation}

\wss{It would be nice to have a blurb here to explain why the subsections below
  cover the requirements.  References to the SRS would be good here.  If a section
  covers tests for input constraints, you should reference the data constraints
  table in the SRS.}

This section covers the tests for project proposal submission and validation, as
specified in FR\_SupervisorProposal, FR\_ProposalOptionalFields, and
FR\_StudentProposal in Section 6 of the SRS. The tests verify successful
supervisor and student proposal submissions, as well as the rejection of
incomplete submissions and acceptance of proposals without optional information.

\begin{enumerate}

\item \textbf{T-FR-01: Submit a Complete Supervisor Proposal}\\

\textbf{Type:} Automatic
					
\textbf{Initial State:} A supervisor is authenticated, has permission to submit
proposals, and the proposal submission period is open.
					
\textbf{Input:} The supervisor completes and submits a proposal with the
following information:
\begin{itemize}
    \item Title: Autonomous Waste Sorting System
    \item Background: Manual waste sorting is inefficient and time-consuming.
    \item Objectives: Develop a system that identifies and sorts recyclable
    materials.
    \item Required Skills: Python, computer vision, and mechanical design.
    \item Deliverables: Working prototype and technical documentation.
    \item Project Impact: Improved recycling efficiency.
    \item Supervisor Capacity: 2 teams
    \item Current Completion: Concept stage
    \item Project Type: Hardware
    \item Industry: Environmental technology
    \item Confidentiality/IP Concerns: None
    \item Additional Terms: No additional restrictions
\end{itemize}
					
\textbf{Output:} The system accepts the submission and stores the proposal with
the ``Proposal submitted'' status making it available for instructors to review.

\textbf{Test Case Derivation:} Derived from \texttt{FR\_SupervisorProposal} of
the SRS. A proposal containing valid information for every field counts as a
valid complete submission. This establishes a normal workflow where incomplete
and invalid submissions can be compared to. Upon acceptance of the input, the
system is expected to accept it and make it available for the instructor to look
through.

\textbf{How test will be performed:} Using Playwright, authenticate as a
supervisor, populate the proposal form with the predefined values, and submit
it. Verify the confirmation and SUBMITTED status. Then, authenticate as an
instructor and confirm that the proposal is accessible with the correct
information.

\item \textbf{T-FR-02: Reject Missing Required Fields}\\

\textbf{Type:} Automatic
					
\textbf{Initial State:} A supervisor is authenticated, has permission to submit
proposals, and the proposal submission period is open.
					
\textbf{Input:} A completed proposal template similar to T-FR-01 with one or
more of the following required fields missing: title, background, course,
objectives, impact, supervisor capacity details, level of current completion,
industry, and information regarding confidentiality and IP concern. The
supervisor will submit. Repeat for each of the following fields.
					
\textbf{Output:} The system rejects the submission, identifies the missing
mandatory field(s), and prevents the proposal from entering the Proposal
submitted state. The instructor will not receive the proposal.

\textbf{Test Case Derivation:} Dervied from \texttt{FR\_ProposalOptionalFields}
of the SRS. Leaving out one mandatory field at a time isolates each validation
constraint and verifies that all required fields are independently enforced.

\textbf{How test will be performed:} Using Playwright, populate the form with a
valid baseline proposal. Remove one mandatory field, attempt submission, and
verify that validation prevents submission. Repeat for each mandatory field. The
test passes if every incomplete proposal is rejected and the corresponding
missing field is identified. Repeated with as many combinations of mandatory
fields as necessary.

\item \textbf{T-FR-03: Accept Missing Optional Fields}\\

\textbf{Type:} Automatic
					
\textbf{Initial State:} A supervisor is authenticated, has permission to submit
proposals, and the proposal submission period is open.
					
\textbf{Input:} A completed proposal template similar to T-FR-01 with all of the
following required fields: title, background, course, objectives, impact,
supervisor capacity details, level of current completion, industry, and
information regarding confidentiality and IP concern. All optional fields will
be left empty. The supervisor will submit. Repeat for each of the following
fields.
					
\textbf{Output:} The system accepts the submission and stores the proposal with
the ``Proposal submitted'' status making it available for instructors to review.

\textbf{Test Case Derivation:} Derived from \texttt{FR\_ProposalOptionalFields}
of the SRS. A proposal containing only mandatory fields represents the minimum
valid submission and verifies that optional information is not incorrectly
enforced as mandatory.

\textbf{How test will be performed:} Using Playwright, complete only the
required proposal fields and submit. Verify that the proposal is accepted
despite the missing optional information and is accessible to the instructor.
The test passes if both conditions are satisfied.

\item \textbf{T-FR-04: Publish a Student-Proposed Project}\\

\textbf{Type:} Automatic
					
\textbf{Initial State:} A student is authenticated and belongs to a team, and
the proposal submission period is open.
					
\textbf{Input:} The student publishes a proposal containing:
\begin{itemize}
    \item Title: Smart Campus Navigation
    \item Problem Statement: Students have difficulty locating classrooms and
    accessible routes across campus.
    \item Proposed Solution/Methodology: Develop a web application providing
    campus navigation and accessible route suggestions.
\end{itemize}
					
\textbf{Output:} The proposal is stored, associated with the submitting team,
and displayed in the separate student-proposed project list with the correct
information.

\textbf{Test Case Derivation:} Derived from \texttt{FR\_StudentProposal}. A
student team providing all specified project information verifies successful
proposal creation and publication in the correct list.

\textbf{How test will be performed:} Use Playwright to authenticate as a student
belonging to a team, complete the student proposal template, and publish it.
Verify the stored information, team association, and visibility in the
student-proposed project list using an account authorized to view it.

\end{enumerate}

\subsubsection{Instructor Project Proposal}

\wss{It would be nice to have a blurb here to explain why the subsections below
  cover the requirements.  References to the SRS would be good here.  If a section
  covers tests for input constraints, you should reference the data constraints
  table in the SRS.}

This section covers FR\_ReviewProposal, including the approval, rejection, and
revision workflows defined in UC\_ReviewProposal. It also accounts for the
updated external supervisor verification rule.

\begin{enumerate}

\item \textbf{T-FR-05: Approve a Supervisor Proposal}\\

\textbf{Control:} Automatic

\textbf{Initial State:} An authenticated instructor has access to an unapproved
proposal titled Autonomous Waste Sorting System submitted by an unverified
external supervisor. The proposal status is ``Proposal submitted''.

\textbf{Input:} The instructor reviews the proposal and approves it.

\textbf{Output:} The project status changes to ``Open'', and the project becomes
visible in the approved catalog to authorized users. The supervisor is notified
and becomes verified.

\textbf{Test Case Derivation:} Derived from \texttt{FR\_ReviewProposal} and
\texttt{FR\_ExternalSupervisorVerification}. Approving an unverified external
supervisor's proposal verifies both the approval transition and the associated
supervisor verification rule.

\textbf{How Test Will Be Performed:} Use Playwright to approve the proposal as
an instructor. Verify the project status, authorized catalog visibility, and
supervisor notification. Then sign in as the external supervisor and verify that
the account is marked verified and can access permitted student profiles. The
test passes if all conditions are satisfied.


\item \textbf{T-FR-06: Request Changes to a Proposal}\\

\textbf{Control:} Automatic

\textbf{Initial State:} An instructor has access to a proposal with the
``Proposal submitted'' status, submitted by an unverified external supervisor.

\textbf{Input:} The instructor requests changes and provides the feedback and
comments asking for clarification on the expected project deliverables and
scope.

\textbf{Output:} The proposal status becomes ``Proposal changes requested''. The
feedback is stored and communicated to the supervisor. The project remains
unpublished, and the supervisor remains unverified.

\textbf{Test Case Derivation:} Derived from \texttt{FR\_ReviewProposal} and
\texttt{FR\_ExternalSupervisorVerification}. A revision request verifies that
feedback does not trigger approval or external supervisor verification.

\textbf{How Test Will Be Performed:} Use Playwright to request changes and
submit the predefined feedback. Verify the status, stored feedback, and
notification. Confirm that the project is excluded from the approved catalog and
that the supervisor remains unverified.


\item \textbf{T-FR-07: Reject a Supervisor Proposal}\\

\textbf{Control:} Automatic

\textbf{Initial State:} An authenticated instructor has access to a proposal
with the ``Proposal submitted'' status, submitted by an unverified external
supervisor.

\textbf{Input:} Reject the proposal with the reason ``The proposed scope does
not meet capstone course expectations.''

\textbf{Output:} The proposal status changes to ``Proposal rejected''. The
rejection reason is preserved, the supervisor is notified, and the project
remains excluded from the approved catalog. The supervisor remains unverified.

\textbf{Test Case Derivation:} Derived from \texttt{FR\_ReviewProposal} and
\texttt{FR\_ExternalSupervisorVerification}. Rejection verifies that an
unsuccessful review does not grant approval-related permissions.

\textbf{How Test Will Be Performed:} Use Playwright to reject the proposal.
Verify the updated status, recorded reason, and notification. Confirm that the
project remains unpublished and the supervisor remains unverified.


\item \textbf{T-FR-08: Revise and Resubmit a Proposal}\\

\textbf{Control:} Automatic

\textbf{Initial State:} A supervisor proposal has the status ``Proposal changes
requested'' and contains instructor feedback asking for clearer project
deliverables.

\textbf{Input:} The supervisor changes the deliverables from ``Working
prototype'' to ``Working prototype, technical documentation, and final
demonstration'', then resubmits.

\textbf{Output:} The revised information is stored, the proposal returns to
``Proposal submitted'', and the instructor can review the updated version.
Previous feedback remains in the review history.

\textbf{Test Case Derivation:} Derived from \texttt{FR\_ReviewProposal} and
\texttt{UC\_ReviewProposal}. Resubmission verifies the transition back to
instructor review without losing earlier feedback.

\textbf{How Test Will Be Performed:} Use Playwright to sign in as the
supervisor, modify the deliverables, and resubmit. Verify the new information
and status. Sign in as the instructor and confirm the revised proposal and
earlier feedback are accessible.

\end{enumerate}

\subsubsection{Project Discovery and Availability}

This section covers \texttt{FR\_ViewApprovedProjects}
and \texttt{FR\_ViewProjectStatus} of the SRS.
Tests verify search and filtering, project
visibility restrictions, and availability
based on reserved supervisor capacity.

\begin{enumerate}

\item \textbf{T-FR-09: Search and Filter Approved Projects}\\

\textbf{Control:} Automatic

\textbf{Initial State:} An authenticated student is enrolled in the course.
Three approved projects are visible to the student:
\begin{itemize}
    \item Project A: Autonomous Waste Sorting System, Environmental Technology, Open.
    \item Project B: Healthcare Scheduling System, Healthcare, Open.
    \item Project C: Automated Assembly Line, Manufacturing, Full.
\end{itemize}

\textbf{Input:} Apply the industry filter ``Environmental Technology'',
availability filter ``Open'', and search keyword ``Autonomous''.

\textbf{Output:} Only Project A is returned. Projects B and C are excluded
because they do not satisfy every criterion.

\textbf{Test Case Derivation:} Derived from \texttt{FR\_ViewApprovedProjects}.
Projects with overlapping and distinct attributes verify that keyword searching
and multiple filters are applied together.

\textbf{How Test Will Be Performed:} Populate the test environment with the
three projects. Use Playwright to apply the filters and keyword. Compare
returned project identifiers with the expected set \{A\}. The test passes if
exactly Project A is returned.



\item \textbf{T-FR-10: Exclude Unapproved or Unauthorized Projects}\\

\textbf{Control:} Automatic

\textbf{Initial State:} The system contains three projects:
\begin{itemize}
    \item Project A: Autonomous Waste Sorting System, approved with status
    ``Open'', submitted by External Supervisor X.
    \item Project B: Drone Inspection Platform, with status ``Proposal
    submitted'', awaiting instructor approval.
    \item Project C: Automated Assembly Line, approved with status ``Full'',
    submitted by External Supervisor Y.
\end{itemize}
A student enrolled in the course, the course instructor, and External
Supervisors X and Y have authenticated accounts. Each external supervisor is
associated only with their own project.

\textbf{Input:} The student searches for Projects A and B. The instructor opens
Project B for review. External Supervisor X searches for Project C, while
External Supervisor Y searches for Project C.

\textbf{Output:} The student can view Project A but cannot view Project B
because it has not been approved. The instructor can access Project B for
review. External Supervisor X cannot view Project C because it belongs to
another external supervisor, while External Supervisor Y can view Project C
because they proposed it.

\textbf{Test Case Derivation:} Derived from \texttt{FR\_ViewApprovedProjects}
and \texttt{FR\_ExternalSupervisorScope} of the SRS. Using projects with
different approval states and ownership verifies that catalog visibility is
restricted by both instructor approval and the user's authorization.

\textbf{How Test Will Be Performed:} Configure the three projects with the
specified statuses and ownership. Using Playwright, authenticate as the student
and search for Projects A and B. Verify that only Project A is displayed.
Authenticate as the instructor and verify that Project B is accessible in the
proposal review interface. Authenticate as External Supervisor X and search for
Project C, verifying that it is excluded from results. Finally, authenticate as
External Supervisor Y and verify that Project C is accessible.

The test passes if each account can access only the project records permitted by
the SRS visibility rules.

\item \textbf{T-FR-11: Display Project Availability and Capacity}\\

\textbf{Control:} Automatic

\textbf{Initial State:} An authenticated student can access two approved
projects:
\begin{itemize}
    \item Project A: Maximum capacity of 2 teams with 1 supervisor slot
    reserved.
    \item Project B: Maximum capacity of 1 team with 1 supervisor slot reserved.
\end{itemize}

\textbf{Input:} The student opens the project catalog and views both projects.

\textbf{Output:} Project A displays ``Open'' with a maximum capacity of 2 teams.
Project B displays ``Full'' with a maximum capacity of 1 team.

\textbf{Test Case Derivation:} Derived from \texttt{FR\_ViewProjectStatus}.
Remaining capacity is calculated as $C_{remaining} = C_{max} - R$ where $R$ is
reserved slot and $C_{max}$ is maximum capacity. Project A has $2-1=1$ available
slot, while Project B has $1-1=0$. These cases verify the correct availability
states.

\textbf{How Test Will Be Performed:} Configure both projects with the specified
capacities and reserved slots. Use Playwright to view their catalog entries and
compare the displayed statuses and maximum capacities with the expected values.


\item \textbf{T-FR-12: Prevent Applications to Full Projects}\\

\textbf{Control:} Automatic

\textbf{Initial State:} An authenticated student belongs to a registered team.
An approved project has a maximum capacity of one team and its only supervisor
slot is reserved, giving it the ``Full'' status.

\textbf{Input:} The student attempts to submit an application to the full
project.

\textbf{Output:} The system rejects the application, informs the student that
the project is unavailable, and does not create an application record.

\textbf{Test Case Derivation:} Derived from \texttt{FR\_ViewProjectStatus} and
\texttt{UC\_SelectProject}. A project with zero remaining capacity verifies that
availability restrictions are enforced by the application workflow rather than
only displayed visually.

\textbf{How Test Will Be Performed:}
Use Playwright to navigate to the full project and attempt to apply through the
normal user interface. Verify that the application is blocked and that no new
application record is created. The test passes if the full project cannot
receive a new application.

\end{enumerate}

\subsubsection{Area of Testing2}

...

\subsection{Tests for Nonfunctional Requirements}

\wss{The nonfunctional requirements for accuracy will likely just reference the
  appropriate functional tests from above.  The test cases should mention
  reporting the relative error for these tests.  Not all projects will
  necessarily have nonfunctional requirements related to accuracy.}

\wss{For some nonfunctional tests, you won't be setting a target threshold for
passing the test, but rather describing the experiment you will do to measure
the quality for different inputs.  For instance, you could measure speed versus
the problem size.  The output of the test isn't pass/fail, but rather a summary
table or graph.}

\wss{Tests related to usability could include conducting a usability test and
  survey.  The survey will be in the Appendix.}

\wss{Static tests, review, inspections, and walkthroughs, will not follow the
format for the tests given below.}

\wss{If you introduce static tests in your plan, you need to provide details.
How will they be done?  In cases like code (or document) walkthroughs, who will
be involved? Be specific.}

\subsubsection{Usability}

..to be completed.

\paragraph{Title for Test}

\begin{enumerate}

\item{test-id1\\}

Type: Functional, Dynamic, Manual, Static etc.
					
Initial State: 
					
Input/Condition: 
					
Output/Result: 
					
How test will be performed: 
					
\item{test-id2\\}

Type: Functional, Dynamic, Manual, Static etc.
					
Initial State: 
					
Input: 
					
Output: 
					
How test will be performed: 

\end{enumerate}

\subsubsection{Installability}

..to be completed.

\subsubsection{Portability}

..to be completed.

...

\subsection{Traceability Between Test Cases and Requirements}

\wss{Provide a table that shows which test cases are supporting which
  requirements.}

\section{Unit Test Description}

\wss{This section should not be filled in until after the MIS (detailed design
  document) has been completed.}

\wss{Reference your MIS (detailed design document) and explain your overall
philosophy for test case selection.}  

\wss{To save space and time, it may be an option to provide less detail in this section.  
For the unit tests you can potentially layout your testing strategy here.  That is, you 
can explain how tests will be selected for each module.  For instance, your test building 
approach could be test cases for each access program, including one test for normal behaviour 
and as many tests as needed for edge cases.  Rather than create the details of the input 
and output here, you could point to the unit testing code.  For this to work, you code 
needs to be well-documented, with meaningful names for all of the tests.}

\subsection{Unit Testing Scope}

\wss{What modules are outside of the scope.  If there are modules that are
  developed by someone else, then you would say here if you aren't planning on
  verifying them.  There may also be modules that are part of your software, but
  have a lower priority for verification than others.  If this is the case,
  explain your rationale for the ranking of module importance.}

\subsection{Tests for Functional Requirements}

\wss{Most of the verification will be through automated unit testing.  If
  appropriate specific modules can be verified by a non-testing based
  technique.  That can also be documented in this section.}

\subsubsection{Module 1}

\wss{Include a blurb here to explain why the subsections below cover the module.
  References to the MIS would be good.  You will want tests from a black box
  perspective and from a white box perspective.  Explain to the reader how the
  tests were selected.}

\begin{enumerate}

\item{test-id1\\}

Type: \wss{Functional, Dynamic, Manual, Automatic, Static etc. Most will
  be automatic}
					
Initial State: 
					
Input: 
					
Output: \wss{The expected result for the given inputs}

Test Case Derivation: \wss{Justify the expected value given in the Output field}

How test will be performed: 
					
\item{test-id2\\}

Type: \wss{Functional, Dynamic, Manual, Automatic, Static etc. Most will
  be automatic}
					
Initial State: 
					
Input: 
					
Output: \wss{The expected result for the given inputs}

Test Case Derivation: \wss{Justify the expected value given in the Output field}

How test will be performed: 

\item{...\\}
    
\end{enumerate}

\subsubsection{Module 2}

...

\subsection{Tests for Nonfunctional Requirements}

\wss{If there is a module that needs to be independently assessed for
  performance, those test cases can go here.  In some projects, planning for
  nonfunctional tests of units will not be that relevant.}

\wss{These tests may involve collecting performance data from previously
  mentioned functional tests.}

\subsubsection{Module ?}
		
\begin{enumerate}

\item{test-id1\\}

Type: \wss{Functional, Dynamic, Manual, Automatic, Static etc. Most will
  be automatic}
					
Initial State: 
					
Input/Condition: 
					
Output/Result: 
					
How test will be performed: 
					
\item{test-id2\\}

Type: Functional, Dynamic, Manual, Static etc.
					
Initial State: 
					
Input: 
					
Output: 
					
How test will be performed: 

\end{enumerate}

\subsubsection{Module ?}

...

\subsection{Traceability Between Test Cases and Modules}

\wss{Provide evidence that all modules have been considered.}
				
\bibliographystyle{plainnat}

\bibliography{../../refs/References}

\newpage

\section{Appendix}

This is where you can place additional information.

\subsection{Symbolic Parameters}

The definition of the test cases will call for SYMBOLIC\_CONSTANTS.
Their values are defined in this section for easy maintenance.

\subsection{Usability Survey Questions?}

\wss{This is a section that would be appropriate for some projects.}

\newpage{}

\section{Appendix --- Generative AI Use Disclosure}

\wss{ChatGPT (version 5) was used to brainstorm ideas for the template for this disclose section.}

\subsection{AI Tools Used}

\wss{Give the name of the AI tool(s) used and the version.}

\subsection{How and Where Used}

\wss{Explain what the tool(s) were used for.  Examples include: brainstorming,
explaining a technical concept, reviewing the document, generating or modifying
text, generating or modifying diagrams, identifying errors or omissions or
inconsistencies.}

\wss{For each of the uses, identify which parts of the document were affected.}

\wss{You do not need to report every interaction, but you should disclose the 
uses that materially contributed to your work}

\subsection{Verification of AI-Generated Content}

\wss{\begin{itemize}
    \item Describe how the team checked AI-generated or AI-assisted material
    before including it in the document.
    \item In particular, verify factual claims, technical assertions,
    requirements, calculations, references, and other information for which an
    AI system may produce incorrect or fabricated results.
    \item \textbf{The team is responsible for the accuracy, completeness,
consistency, and appropriateness of the submitted document, regardless of
whether material was generated by a human or an AI system.}
\end{itemize}}

\newpage{}

\section{Appendix --- Reflection}

\wss{This section is not required for CAS 741}

The information in this section will be used to evaluate the team members on the
graduate attribute of Lifelong Learning.

\input{../Reflection.text}

\begin{enumerate}
  \item What went well while writing this deliverable? 
  \item What pain points did you experience during this deliverable, and how
    did you resolve them?
  \item What knowledge and skills will the team collectively need to acquire to
  successfully complete the verification and validation of your project?
  Examples of possible knowledge and skills include dynamic testing knowledge,
  static testing knowledge, specific tool usage, Valgrind etc.  You should look to
  identify at least one item for each team member.
  \item For each of the knowledge areas and skills identified in the previous
  question, what are at least two approaches to acquiring the knowledge or
  mastering the skill?  Of the identified approaches, which will each team
  member pursue, and why did they make this choice?
\end{enumerate}

\end{document}